\begin{comment}
El proyecto presentado tiene como objetivo el diseño y fabricación de una impresora 3D de modelado por deposición fundida (FDM) optimizada para el uso de filamento PLA en entornos estudiantiles. El enfoque principal es proporcionar a los alumnos una plataforma de experimentación práctica y asequible que permita comprender a fondo la tecnología de Control Numérico por Computadora (CNC) y la manufactura aditiva.
 
Mediante el uso de herramientas de diseño asistido por computadora (CAD) como SolidWorks y plataformas de desarrollo de hardware abierto como Arduino, se ha desarrollado un sistema capaz de ser fabricado, en gran medida, por otra impresora 3D. Este concepto de "máquina autorreplicante" no solo reduce costos, sino que destaca la importancia del ensamblaje, la calibración y el ajuste mecánico como factores determinantes en la precisión final del equipo.

El desarrollo de esta impresora permite al estudiante involucrarse directamente en procesos de diseño, manufactura y calibración de sistemas de manufactura aditiva. Al enfrentarse al desafío de diseñar y fabricar la máquina desde sus componentes base, el estudiante desarrolla habilidades técnicas críticas y una capacidad de resolución de problemas que va más allá del uso de herramientas comerciales cerradas.
\end{comment}

La manufactura aditiva mediante Modelado por Deposición Fundida (FDM) representa hoy una de las tecnologías de mayor impacto en entornos educativos y de prototipado, debido a su bajo costo de implementación y a la flexibilidad que ofrece para la fabricación rápida de componentes funcionales. En particular, las impresoras 3D basadas en arquitecturas abiertas han permitido que estudiantes y desarrolladores puedan intervenir directamente en aspectos mecánicos, electrónicos y de control, favoreciendo el aprendizaje práctico de sistemas de manufactura automatizada.

El presente proyecto tiene como objetivo el diseño y fabricación de una impresora 3D tipo FDM orientada a aplicaciones académicas, utilizando filamento PLA como material principal de manufactura. El sistema fue desarrollado a partir de una arquitectura de hardware abierto inspirada en la filosofía RepRap. Para ello, 
se partió de la rehabilitación de un equipo institucional fuera de servicio, el cual no solo sirvió como referencia para identificar las limitaciones de las configuraciones comerciales, sino que se utilizó como herramienta de apoyo para fabricar los múltiples componentes estructurales de la nueva impresora, facilitando así futuras modificaciones mecánicas y electrónicas.

Para el desarrollo del equipo se emplearon herramientas de diseño asistido por computadora (CAD), específicamente SolidWorks, así como plataformas de control electrónico basadas en Arduino y firmware Marlin. Adicionalmente, se realizó un rediseño del módulo de extrusión y del sistema de enfriamiento de capa con el propósito de mejorar la estabilidad térmica y la calidad superficial de las piezas.Este desarrollo respondió a criterios de adaptación al contexto académico, incluyendo consideraciones de costo, disponibilidad de materiales y mejora del comportamiento térmico.

El desarrollo de esta plataforma permite al estudiante involucrarse directamente en el ciclo completo de un sistema de manufactura aditiva: desde el diseño mecánico y la integración electrónica hasta la calibración, el diagnóstico de fallas (como la absorción de humedad en el filamento) y el control numérico. De esta manera, el proyecto no solo representa una alternativa de bajo costo frente a equipos comerciales cerrados, sino también una herramienta didáctica capaz de fortalecer competencias técnicas relacionadas con sistemas embebidos, automatización y manufactura digital.

Finalmente, el presente documento se organiza de la siguiente manera: en el Capítulo 2 se expone el planteamiento del problema; en el Capítulo 3 se 
presentan los antecedentes y la justificación del proyecto; en el Capítulo 4 se definen los objetivos generales, específicos y las metas planteadas. 
Posteriormente, el Capítulo 5 desarrolla el marco teórico. En el Capítulo 6 se describe la metodología, que incluye la rehabilitación del equipo, el diseño del módulo de extrusión y los procesos de calibración. El Capítulo 7 presenta los resultados obtenidos y el Capítulo 8 expone la discusión de dichos resultados, contrastándolos con la literatura y reconociendo las 
limitaciones del trabajo. Finalmente, el Capítulo 9 presenta las conclusiones del proyecto, y el Capítulo 10 incluye las referencias bibliográficas que sustentan la investigación.
